% K31-30, JC-003 y P01 del inventario REV05.
% Residencia: después de la derivación del registro K y de la incidencia Witt;
% antes de utilizar la lectura ponderada para la recuperación del registro.
% Propietarios: output/IMPLEMENTACION_K_SEMILLAS_INCIDENCIA_20260918/
% WeightedIncidenceRecovery.lean (íntegro) y
% COMPOSICION_GENERATIVA_Y_RECUPERACION_K.tex:257--306.
% Codificador: PAQUETE_ARTICULO_I_INTEGRACION_ESTADO_CAMPOS_20260919/
% lean/biblioteca/CoxeterNeighbor.lean:322--343;
% lean/biblioteca/SectorIncidenceData.lean:27--33.
% Estatuto: RESULTADO_RECUPERADO, explicitación íntegra de las matrices que
% la exposición anterior remitía al código. La prueba Lean es la ya existente.

\section[Récupération par incidences hexadiques]{Récupération intégrale du registre au moyen d'incidences hexadiques pondérées}
\label{rh:sec:recuperacion}

Le registre dodécaphasique engendré conserve douze composantes ordonnées et
une origine de phase. L'incidence de Paley--Witt fournit une lecture
supplémentaire de ce registre : chaque support hexadique rassemble six de ses
composantes au moyen d'une somme pondérée par l'incidence elle-même. La
reconstruction développée ici reçoit les valeurs de ces sommes
du même registre. Son objet est de prouver que les lectures conservent
toutes les composantes et d'exhiber l'opérateur inverse à coefficients
rationnels explicites.

\subsection{Codeur d'incidence et douze témoins explicites}

Dans le système de coordonnées d'incidence déjà fixé, soit
\[
A_W=\begin{pmatrix}
0&1&1&1&1&1\\1&0&1&2&2&1\\1&1&0&1&2&2\\
1&2&1&0&1&2\\1&2&2&1&0&1\\1&1&2&2&1&0
\end{pmatrix}\in M_6(\mathbb F_3).
\]
Pour $w\in\mathbb F_3^6$, le codeur est
$\operatorname{Enc}(w)=(w,wA_W)\in\mathbb F_3^{12}$ et son support
est
\[
H(w)=\{j\in\{1,\ldots,12\}:\operatorname{Enc}(w)_j\ne0\}.
\]
Le choix d'une position utilise des indices de un à douze ; les
implémentations à indices commençant à zéro sont transportées au moyen de
$j\mapsto j+1$.

Les douze témoins suivants appartiennent au domaine du codeur :
\[
\begin{array}{c|c@{\qquad}c|c}
i&w^{(i)}&i&w^{(i)}\\\hline
1&112210&7&122020\\
2&122101&8&112002\\
3&121200&9&111001\\
4&111100&10&110122\\
5&121012&11&101221\\
6&111020&12&011111.
\end{array}
\]
Définissons $H_i=H(w^{(i)})$ et
$B_{ij}=\mathbf1_{\{j\in H_i\}}$. L'évaluation du codeur donne
\begingroup\small\setlength{\arraycolsep}{3pt}
\setcounter{MaxMatrixCols}{12}
\begin{equation}
B=\begin{pmatrix}
1&1&1&1&1&0&0&0&0&0&0&1\\
1&1&1&1&0&1&0&0&0&0&1&0\\
1&1&1&1&0&0&1&0&1&0&0&0\\
1&1&1&1&0&0&0&1&0&1&0&0\\
1&1&1&0&1&1&0&0&0&1&0&0\\
1&1&1&0&1&0&1&0&0&0&1&0\\
1&1&1&0&1&0&0&1&1&0&0&0\\
1&1&1&0&0&1&1&1&0&0&0&0\\
1&1&1&0&0&1&0&0&1&0&0&1\\
1&1&0&1&1&1&0&0&1&0&0&0\\
1&0&1&1&1&1&0&1&0&0&0&0\\
0&1&1&1&1&1&1&0&0&0&0&0
\end{pmatrix}.
\label{rh:eq:incidencia}
\end{equation}
\endgroup
Chaque ligne contient six entrées égales à un. Par conséquent, chaque
$H_i$ est une hexade de l'image du codeur, et la matrice $B$ est
construite à partir de ces supports.

\subsection{Inverse rationnel et unicité de récupération}

Soit la matrice entière
\begingroup\small\setlength{\arraycolsep}{3pt}
\setcounter{MaxMatrixCols}{12}
\begin{equation}
M=\begin{pmatrix}
7&1&8&-7&7&-1&-4&8&-7&3&3&-15\\
7&7&-4&-1&1&-7&8&8&-7&3&-15&3\\
1&7&8&-7&7&-7&8&-4&-1&-15&3&3\\
1&1&2&5&-5&-1&-4&-4&-1&3&3&3\\
1&-5&-4&-1&1&5&2&-4&-1&3&3&3\\
-5&1&-4&-1&1&-1&-4&2&5&3&3&3\\
-5&-11&2&5&-5&11&-10&2&5&3&3&3\\
-5&-5&-10&11&-11&5&2&2&5&3&3&3\\
-11&-5&2&5&-5&5&2&-10&11&3&3&3\\
-11&-11&-4&17&1&11&-10&-10&11&3&3&3\\
-11&1&-10&11&-11&17&-4&-10&11&3&3&3\\
1&-11&-10&11&-11&11&-10&-4&17&3&3&3
\end{pmatrix}.
\label{rh:eq:inversa-entera}
\end{equation}
\endgroup

\begin{theorem}[Récupération exacte à partir de douze lectures hexadiques]
\label{rh:thm:doce}
Pour $v\in\mathbb Q^{12}$, soient
\[
\operatorname{Read}(v)=Bv,
\qquad
\operatorname{Rec}(y)=\frac1{18}My.
\]
On vérifie
\[
\operatorname{Rec}(\operatorname{Read}(v))=v,
\qquad
\operatorname{Read}(\operatorname{Rec}(y))=y.
\]
Par conséquent, tout $y\in\mathbb Q^{12}$ admet un unique registre
$v\in\mathbb Q^{12}$ avec $\operatorname{Read}(v)=y$.
\end{theorem}
\begin{proof}
Les multiplications entières des matrices
\eqref{rh:eq:incidencia} et \eqref{rh:eq:inversa-entera} donnent
$MB=BM=18I_{12}$. Chaque entrée diagonale résulte d'un produit
scalaire égal à dix-huit et chaque entrée hors diagonale résulte
d'un produit scalaire nul. En divisant par dix-huit, on obtient
les deux identités inverses.

En particulier,
$\operatorname{Rec}(\operatorname{Read}(v))=(MB/18)v=v$ et
$\operatorname{Read}(\operatorname{Rec}(y))=(BM/18)y=y$.
La seconde identité prouve l'existence, au moyen de
$v=\operatorname{Rec}(y)$. Si $Bu=Bv$, appliquer $M/18$ aux deux
membres donne $u=v$, ce qui prouve l'unicité.
\end{proof}

La vérification de $B$ à partir du codeur et les deux multiplications
exactes sont formalisées dans les théorèmes
\texttt{incidence\_from\_encoder}, \texttt{inverse\_left} et
\texttt{inverse\_right}. Les identités de composition correspondent
à \texttt{recover\_read} et \texttt{read\_recover}.
Leur contenu mathématique se trouve intégralement dans les matrices et
dans la preuve précédente ; le certificat formel consigne la vérification
de ces mêmes opérations.

\begin{corollary}[Séparation par toutes les incidences pondérées]
Si $u,v\in\mathbb Q^{12}$ satisfont
\[
\sum_{j\in H(w)}u_j=\sum_{j\in H(w)}v_j
\]
pour tout $w\in\mathbb F_3^6$ dont le support codé comporte six
positions, alors $u=v$.
\end{corollary}
\begin{proof}
L'hypothèse s'applique en particulier aux douze témoins exposés.
Par la définition de $B$, leurs douze égalités constituent $Bu=Bv$.
Le théorème~\ref{rh:thm:doce} permet de récupérer les deux entrées et
conclut $u=v$.
\end{proof}

\begin{example}[Récupération d'une composante périphérique]
Pour le vecteur canonique $v=e_{12}$, les seules lectures non nulles sont
la première et la neuvième :
\[
y=Be_{12}=(1,0,0,0,0,0,0,0,1,0,0,0)^{\mathsf T}.
\]
La somme des première et neuvième colonnes de $M$ est $18e_{12}$.
Par conséquent, $My/18=e_{12}$. Les deux sommes hexadiques distinguent
cette position dans la lecture conjointe parce que les autres composantes
s'annulent exactement dans l'opérateur de récupération.
\end{example}

\subsection{Opérateur de récupération symétrique sur les cent trente-deux hexades}

Le système de Steiner $S(5,6,12)$ déjà produit par l'incidence
contient $132$ hexades. Notons leur ensemble $\mathcal H$ et
définissons
\[
\mathcal I:\mathbb Q^{12}\longrightarrow\mathbb Q^{\mathcal H},
\qquad
(\mathcal I v)_H=\sum_{j\in H}v_j.
\]
Ce lecteur présente une redondance : il conserve les douze composantes au moyen de
$132$ sommes, tandis que le théorème précédent fournit une base de
douze de ces sommes.

\begin{theorem}[Inverse sur l'image du lecteur complet]
\label{rh:thm:steiner}
On a
\[
\mathcal I^{\mathsf T}\mathcal I=36I_{12}+30J_{12},
\]
où $J_{12}$ a toutes ses entrées égales à un. Pour
$y=\mathcal I v$, le registre est récupéré au moyen de
\begin{equation}
396v_i=11\sum_{H\ni i}y_H-5\sum_{H\in\mathcal H}y_H,
\qquad i=1,\ldots,12.
\label{rh:eq:recuperador}
\end{equation}
\end{theorem}
\begin{proof}
La propriété $S(5,6,12)$ implique que chaque position appartient à
$\binom{11}{4}/\binom{5}{4}=66$ hexades. Chaque paire de positions
distinctes appartient à $\binom{10}{3}/\binom{4}{3}=30$ hexades.
Par conséquent, le produit de la transposée de la matrice d'incidence par
elle-même a pour diagonale $66$ et pour entrées hors diagonale $30$, ce qui
est $36I+30J$.

Soit $S=\sum_jv_j$. Les deux comptages précédents donnent
\[
\sum_{H\ni i}y_H=66v_i+30\sum_{j\ne i}v_j=36v_i+30S,
\qquad
\sum_Hy_H=66S.
\]
Multiplier la première égalité par onze et soustraire cinq fois la
seconde élimine $S$ et produit $396v_i$. Cela démontre
\eqref{rh:eq:recuperador} pour toute lecture appartenant à
l'image de $\mathcal I$.
\end{proof}

\begin{corollary}[Composantes uniforme et transversale]
Sur $\mathbb R^{12}$, le lecteur complet a pour valeur singulière
$\sqrt{396}$ dans la direction uniforme et pour valeur singulière $6$ dans son
complément orthogonal. Son opérateur inverse sur l'image a pour norme
opératorielle $1/6$.
\end{corollary}
\begin{proof}
La matrice $J$ agit par douze sur $\mathbb R\mathbf1$ et par zéro
sur $\mathbf1^\perp$. Les valeurs propres de $36I+30J$ sont, par
conséquent, $396$ et $36$. Les valeurs singulières sont leurs racines
positives ; la plus grande valeur singulière de l'opérateur inverse est l'inverse
de la plus petite d'entre elles.
\end{proof}

\subsection{Composition avec le registre engendré}

En appliquant ces identités au registre entier dodécaphasique $K$, déjà
produit avec son ordre et son origine, les lectures $BK$ ou $\mathcal I K$
permettent de récupérer exactement ses douze composantes. L'interprétation
positionnelle ultérieure prend ce même registre en base mille ; les
opérateurs de Hadamard, les différences et les lecteurs hexadiques sont
des réalisations récupérables d'une provenance commune.

Le système de douze coordonnées rationnelles admet toute entrée
$y\in\mathbb Q^{12}$ et la reconstruit comme un registre rationnel
unique. Le sous-domaine correspondant aux registres HMT engendrés
conserve, en outre, ses conditions d'intégralité, de largeur des
triades, d'origine, de sélection et de frontière. Pour le système redondant de
$132$ lectures, l'inverse agit sur l'image de $\mathcal I$.
Ces précisions typent le contenu récupéré sans modifier la
génération antécédente du registre.
